% Fazzle — Apresentação Beamer (CEFET-MG Campus Timóteo)
% Compilação: pdflatex (2 passes) — ex.: ./compile.sh fazzle
% Contraste: texto claro sobre VerdeEscuro/VerdeFloresta; texto VerdeEscuro
%           sobre AmareloFauna/VerdeCitrico (meta WCAG 4.5:1 — validar em WebAIM se necessário).

\documentclass[aspectratio=169,11pt]{beamer}

\usepackage[T1]{fontenc}
\usepackage[utf8]{inputenc}
\usepackage[brazil]{babel}
\usepackage{lmodern}
\usepackage{microtype}

\usepackage[table]{xcolor}
\definecolor{VerdeFloresta}{HTML}{519E5D}
\definecolor{VerdeEscuro}{HTML}{1E4A28}
\definecolor{AmareloFauna}{HTML}{F2B705}
\definecolor{VerdeCitrico}{HTML}{94C12D}
\definecolor{LaranjaAlerta}{HTML}{E87C1E}

\usepackage{tikz}
\usetikzlibrary{calc,positioning,backgrounds}
\usepackage[most]{tcolorbox}
\tcbuselibrary{skins}

\usepackage{booktabs}

% --- Beamer: identidade Fazzle / CEFET ---
\usetheme{Madrid}
\usecolortheme{default}
\setbeamercolor{structure}{fg=VerdeEscuro}
\setbeamercolor{frametitle}{bg=VerdeEscuro,fg=white}
\setbeamercolor{title}{bg=VerdeEscuro,fg=white}
\setbeamercolor{subtitle}{fg=AmareloFauna}
\setbeamercolor{author in head/foot}{bg=VerdeFloresta,fg=white}
\setbeamercolor{title in head/foot}{bg=VerdeEscuro,fg=white}
\setbeamercolor{date in head/foot}{bg=VerdeFloresta,fg=white}
\setbeamercolor{block title}{bg=VerdeFloresta,fg=white}
\setbeamercolor{block body}{bg=black!3,fg=black}
\setbeamercolor{block title alerted}{bg=LaranjaAlerta,fg=white}
\setbeamercolor{block body alerted}{bg=black!5,fg=black}
\setbeamertemplate{navigation symbols}{}

% Caixas “marca-texto” e CTA (contraste preferencial)
\newtcolorbox{fazzlehl}[1][]{%
  enhanced,
  colback=AmareloFauna!88,
  colframe=VerdeEscuro,
  coltitle=VerdeEscuro,
  coltext=VerdeEscuro,
  fonttitle=\bfseries,
  arc=3mm,
  boxrule=0.8pt,
  left=8pt,
  right=8pt,
  top=6pt,
  bottom=6pt,
  #1}
\newtcolorbox{fazzlehlcitrico}[1][]{%
  enhanced,
  colback=VerdeCitrico!35,
  colframe=VerdeEscuro,
  coltitle=VerdeEscuro,
  coltext=VerdeEscuro,
  fonttitle=\bfseries,
  arc=3mm,
  boxrule=0.8pt,
  left=8pt,
  right=8pt,
  top=6pt,
  bottom=6pt,
  #1}
\newtcolorbox{fazzlecta}[1][]{%
  enhanced,
  colback=LaranjaAlerta,
  colframe=VerdeEscuro,
  colupper=white,
  fontupper=\bfseries\large,
  arc=4mm,
  boxrule=1pt,
  left=10pt,
  right=10pt,
  top=8pt,
  bottom=8pt,
  #1}

\newcommand{\tituloimpacto}[1]{\textcolor{LaranjaAlerta}{\textbf{#1}}}

% Fundo decorativo (círculos) — título
\newcommand{\FazzleTituloBG}{%
  \begin{tikzpicture}[remember picture,overlay]
    \fill[VerdeFloresta,opacity=0.18] ($(current page.north west)+(1.2,-1.2)$) circle (2.8cm);
    \fill[AmareloFauna,opacity=0.14] ($(current page.north east)+(-2,-2.5)$) circle (3.4cm);
    \fill[VerdeCitrico,opacity=0.12] ($(current page.south west)+(3,2)$) circle (2.2cm);
    \fill[LaranjaAlerta,opacity=0.10] ($(current page.south east)+(-2.5,1.5)$) circle (2.5cm);
  \end{tikzpicture}%
}

% Três nós circulares (ecossistema)
\newcommand{\FazzleTresCirculos}[3]{%
  \begin{center}
  \begin{tikzpicture}[
      node distance=1.6cm,
      every node/.style={align=center,inner sep=6pt}]
    \node[circle,draw=VerdeEscuro,line width=1pt,fill=VerdeFloresta!25,minimum size=2.9cm,text width=2.4cm] (A) {\small\bfseries #1};
    \node[circle,draw=VerdeEscuro,line width=1pt,fill=AmareloFauna!45,minimum size=2.9cm,text width=2.4cm,right=of A] (B) {\small\bfseries #2};
    \node[circle,draw=VerdeEscuro,line width=1pt,fill=VerdeCitrico!40,minimum size=2.9cm,text width=2.4cm,right=of B] (C) {\small\bfseries #3};
    \draw[VerdeEscuro!60,thick,-stealth] (A) -- (B);
    \draw[VerdeEscuro!60,thick,-stealth] (B) -- (C);
  \end{tikzpicture}
  \end{center}}

\title[Fazzle]{Fazzle: Exploradores da Fauna Brasileira}
\subtitle{Brincar, sentir e proteger. Quebra-cabeça multissensorial para educação ambiental e inclusão.}
\author{Projeto Fazzle}
\institute{CEFET-MG --- Campus Timóteo}
\date{\today}

\begin{document}

% --- Capa ---
\begin{frame}[plain]
  \FazzleTituloBG
  \vspace*{0.5cm}
  \begin{center}
    {\usebeamercolor[fg]{title}\usebeamerfont{title}\inserttitle\par}
    \vspace{0.6em}
    {\usebeamercolor[fg]{subtitle}\usebeamerfont{subtitle}\insertsubtitle\par}
    \vspace{1.2em}
    \begin{fazzlehl}[title=Proposta em uma linha]
      Conectamos \textbf{tato}, \textbf{áudio} e \textbf{narrativa ecológica} com tecnologia de proximidade sem telas para a criança explorar a fauna brasileira com autonomia.
    \end{fazzlehl}
    \vspace{1em}
    {\small\insertinstitute\par}
  \end{center}
\end{frame}

% --- Problema e oportunidade ---
\begin{frame}{Lacuna na educação inclusiva}
  \begin{columns}[T,onlytextwidth]
    \begin{column}{0.48\textwidth}
      \begin{block}{O problema}
        \begin{itemize}\setlength\itemsep{0.35em}
          \item Brinquedos táteis focam em \textbf{formas básicas}, com pouca profundidade narrativa.
          \item Jogos de conscientização ecológica \textbf{dependem demais da visão}.
        \end{itemize}
      \end{block}
    \end{column}
    \begin{column}{0.48\textwidth}
      \begin{fazzlehlcitrico}[title=Oceano azul]
        Pais e educadores buscam soluções \textbf{STEAM}, \textbf{acessíveis} e \textbf{sustentáveis}. O mercado global de brinquedos educacionais cresce de forma acelerada até 2032.
      \end{fazzlehlcitrico}
    \end{column}
  \end{columns}
  \vspace{0.6em}
  \begin{fazzlecta}
    Atuamos na interseção: inclusão + ambiental + tecnologia assistiva --- ainda pouco explorada no varejo educacional.
  \end{fazzlecta}
\end{frame}

% --- Proposta de valor (neuromotor + PCD visual) ---
\begin{frame}{Proposta de valor}
  \begin{fazzlehl}[title=Para quem importa]
    Estimulamos \textbf{memória de trabalho}, \textbf{coordenação fino-motora} e \textbf{associação semiótica} por meio do encaixe magnético, texturas e narrativa sonora --- com foco em crianças com \textbf{deficiência visual} e baixa visão, sem excluir o público amplo.
  \end{fazzlehl}
  \vspace{0.5em}
  \begin{itemize}\setlength\itemsep{0.4em}
    \item \textbf{Alto contraste} e orientação tátil reduzem carga cognitiva na exploração.
    \item \textbf{Áudio contextual} reforça vocabulário, habitat e conservação.
    \item \textbf{Desafio progressivo} (``just right challenge'') mantém engajamento terapêutico e lúdico.
  \end{itemize}
\end{frame}

% --- Ecossistema do produto ---
\begin{frame}{Ecossistema Fazzle: tato, proximidade e áudio}
  \FazzleTresCirculos{Ação física}{Identificação por proximidade}{Recompensa auditiva}
  \vspace{0.4em}
  \begin{itemize}\setlength\itemsep{0.35em}
    \item A criança \textbf{tateia} a peça texturizada e \textbf{encaixa} na base magnética.
    \item Responsáveis \textbf{aproximam o smartphone} da peça (tags \textbf{NFC} --- comunicação por proximidade, sem bateria na tag).
    \item O aplicativo \textbf{narra histórias} e reproduz sons reais da fauna --- \tituloimpacto{sem telas obrigatórias para a brincadeira}.
  \end{itemize}
\end{frame}

% --- Design universal ---
\begin{frame}{Design universal na prática}
  \begin{columns}[T,onlytextwidth]
    \begin{column}{0.32\textwidth}
      \begin{tcolorbox}[colback=VerdeFloresta!20,colframe=VerdeEscuro,arc=2mm,title=\textcolor{VerdeEscuro}{\textbf{Orientação tátil}}]
        \small Entalhes indicam o topo; integração com \textbf{Braille} nas peças.
      \end{tcolorbox}
    \end{column}
    \begin{column}{0.32\textwidth}
      \begin{tcolorbox}[colback=AmareloFauna!50,colframe=VerdeEscuro,arc=2mm,title=\textcolor{VerdeEscuro}{\textbf{Alto contraste}}]
        \small Cores vibrantes sobre fundo claro para baixa visão.
      \end{tcolorbox}
    \end{column}
    \begin{column}{0.32\textwidth}
      \begin{tcolorbox}[colback=VerdeCitrico!30,colframe=VerdeEscuro,arc=2mm,title=\textcolor{VerdeEscuro}{\textbf{Base magnética}}]
        \small O ``clique'' do ímã dá \textbf{previsibilidade} e evita quedas durante a exploração.
      \end{tcolorbox}
    \end{column}
  \end{columns}
\end{frame}

% --- Impacto educacional ---
\begin{frame}{Impacto educacional}
  \begin{block}{Dinâmica principal}
    Associação guiada: agrupar herbívoros, encaixar o animal ao \textbf{esqueleto em relevo} e consolidar vocabulário ecológico.
  \end{block}
  \begin{fazzlehl}[title=Tecnologia de proximidade]
    Ao aproximar o smartphone do quebra-cabeça montado, a criança ouve o \textbf{som real do animal}, histórias de habitat e alertas sobre riscos de extinção --- conteúdo curado e acessível.
  \end{fazzlehl}
\end{frame}

% --- Sustentabilidade ---
\begin{frame}{Materialidade e mensagem}
  \begin{itemize}\setlength\itemsep{0.4em}
    \item Superamos papelão frágil e plásticos descartáveis: \textbf{cortiça natural}, \textbf{bioplásticos} e \textbf{madeiras certificadas FSC}.
    \item A base magnética é \textbf{removível}: o desafio evolui com a criança (escalonamento de dificuldade).
  \end{itemize}
  \vspace{0.5em}
  \begin{center}
    \begin{tabular}{@{}lcc@{}}
      \toprule
      & \textbf{Fazzle} & \textbf{Tradicional} \\
      \midrule
      Sensorial & Tátil + áudio & Predominantemente visual \\
      Sustentabilidade & Materiais de baixo impacto & Papelão / plástico barato \\
      \bottomrule
    \end{tabular}
  \end{center}
\end{frame}

% --- Segmentação ---
\begin{frame}{Segmentação}
  \begin{fazzlecta}
    Crianças de \textbf{5 a 10 anos} com deficiência visual (e baixa visão) + \textbf{responsáveis}, educadores e terapeutas que buscam brinquedos com evidência de engajamento motor e narrativo.
  \end{fazzlecta}
  \vspace{0.6em}
  \begin{itemize}\setlength\itemsep{0.35em}
    \item Mensagem clara para \textbf{pais}: segurança, STEAM e inclusão.
    \item Canal técnico para \textbf{escolas e clínicas}: kits e protocolos de uso.
  \end{itemize}
\end{frame}

% --- Parcerias / atividades / recursos ---
\begin{frame}{Motor do projeto}
  \begin{columns}[T,onlytextwidth]
    \begin{column}{0.34\textwidth}
      \begin{tcolorbox}[colback=VerdeEscuro!8,colframe=VerdeFloresta,arc=2mm,title=\textcolor{VerdeEscuro}{\textbf{Parcerias}}]
        \small Ecossistema de suprimentos: informática, química e embalagens. \textbf{Marketing de influência} para prova social e alcance orgânico.
      \end{tcolorbox}
    \end{column}
    \begin{column}{0.32\textwidth}
      \begin{tcolorbox}[colback=AmareloFauna!40,colframe=VerdeEscuro,arc=2mm,title=\textcolor{VerdeEscuro}{\textbf{Atividades-chave}}]
        \small Pesquisa e desenvolvimento de \textbf{hardware e software} para acessibilidade, UX infantil e educação inclusiva.
      \end{tcolorbox}
    \end{column}
    \begin{column}{0.32\textwidth}
      \begin{tcolorbox}[colback=VerdeCitrico!28,colframe=VerdeEscuro,arc=2mm,title=\textcolor{VerdeEscuro}{\textbf{Recursos}}]
        \small Expertise em \textbf{game design} ambiental, laboratório de prototipagem e infraestrutura computacional e física no campus.
      \end{tcolorbox}
    \end{column}
  \end{columns}
\end{frame}

% --- Custos ---
\begin{frame}{Estrutura de custos (frentes de investimento)}
  \begin{itemize}\setlength\itemsep{0.45em}
    \item \textbf{Aplicativo} multiplataforma (\textbf{Android / iOS}): narrativa, áudio, acessibilidade e atualização de conteúdos.
    \item \textbf{Integração por proximidade} (NFC) e testes de interoperabilidade com dispositivos-alvo.
    \item \textbf{Magnetismo} industrial: ímãs seguros, fixação e durabilidade.
    \item \textbf{Bioplástico} e biomateriais: matéria-prima, moldagem e certificações.
  \end{itemize}
  \vspace{0.4em}
  \begin{fazzlehlcitrico}[title=Gestão]
    Priorizamos lote-piloto, homologação de fornecedores e custo total de propriedade (não apenas preço unitário).
  \end{fazzlehlcitrico}
\end{frame}

% --- Canais, relacionamento, receita ---
\begin{frame}{Canais, relacionamento e receita}
  \begin{columns}[T,onlytextwidth]
    \begin{column}{0.5\textwidth}
      \begin{block}{Canais e relacionamento}
        \begin{itemize}\setlength\itemsep{0.3em}
          \item \textbf{B2C}: e-commerce, redes sociais e campanhas de \textbf{AdTech} segmentadas.
          \item Conteúdo educativo contínuo para \textbf{comunidade} e suporte pós-venda.
        \end{itemize}
      \end{block}
    \end{column}
    \begin{column}{0.48\textwidth}
      \begin{fazzlecta}
        Receita principal: \textbf{venda direta transacional} digital (checkout próprio ou marketplaces).
      \end{fazzlecta}
    \end{column}
  \end{columns}
\end{frame}

% --- Go-to-market ---
\begin{frame}{Como chegamos ao mercado}
  \begin{itemize}\setlength\itemsep{0.4em}
    \item \textbf{Varejo B2C}: preço-alvo de referência \textbf{R\$50--R\$80} (e-commerce direto ao consumidor).
    \item \textbf{B2B institucional}: escolas inclusivas, clínicas de terapia ocupacional e institutos para pessoas com deficiência visual.
    \item \textbf{Escalabilidade}: kits colecionáveis por biomas e espécies da fauna brasileira.
  \end{itemize}
\end{frame}

% --- Fecho ---
\begin{frame}[plain]
  \FazzleTituloBG
  \vspace*{1.2cm}
  \begin{center}
    {\usebeamercolor[fg]{title}\usebeamerfont{title}Próximos passos\par}
    \vspace{1em}
    \begin{fazzlehl}[title=Convite]
      Agendamos demonstração do protótipo, parceria de conteúdo científico e \textbf{piloto institucional} com o CEFET-MG Campus Timóteo.
    \end{fazzlehl}
    \vspace{1.2em}
    {\large\textcolor{VerdeEscuro}{\textbf{contato@cefetmg.br}} \quad\textcolor{black!60}{\small(defina o e-mail oficial do projeto)}\par}
    \vspace{0.5em}
    {\small\insertinstitute\ --- Projeto Fazzle\par}
  \end{center}
\end{frame}

\end{document}
